% !TEX TS-program = xelatex
% !TeX spellcheck = ru_RU
% Файл можно включить в диплом командой % !TEX TS-program = xelatex
% !TeX spellcheck = ru_RU
% Файл можно включить в диплом командой % !TEX TS-program = xelatex
% !TeX spellcheck = ru_RU
% Файл можно включить в диплом командой % !TEX TS-program = xelatex
% !TeX spellcheck = ru_RU
% Файл можно включить в диплом командой \input{abstract}.
% При самостоятельной компиляции используется оформление предпросмотра.
\ifdefined\maketitle
    \let\abstractdocumentend\relax
\else
    \documentclass[12pt,a4paper]{article}
    \usepackage{fontspec}
    \setmainfont{cmunrm.otf}[BoldFont=cmunbx.otf,ItalicFont=cmunti.otf,BoldItalicFont=cmunbi.otf]
    \usepackage{polyglossia}
    \setdefaultlanguage{russian}
    \usepackage[top=20mm,bottom=20mm,left=30mm,right=15mm]{geometry}
    \usepackage{setspace}
    \onehalfspacing
    \setlength{\emergencystretch}{2em}
    \def\abstractdocumentend{\end{document}}
    \begin{document}
\fi

\section*{Аннотация}

Работа посвящена распределению ограниченной памяти между дообучением модели
и калибровкой порога сетевого детектора в условиях возвратного дрейфа трафика. Под возвратным дрейфом понимается повторное появление
ранее наблюдавшихся режимов сетевой активности после периода их отсутствия.
Например, дневная активность пользователей в сети организации
сменяется снижением нагрузки ночью. После адаптации к ночному
трафику детектор может ошибочно выдавать тревоги при возвращении
обычной дневной активности на следующее утро.
При непрерывном обновлении детектора такие ошибки могут возникать как
вследствие ухудшения распознавания прежних данных, так и из-за изменения
численных оценок аномальности, которым перестаёт соответствовать установленный
порог.

Для сохранения информации о прошлом используются исторические наблюдения.
Часть из них может служить повторному обучению модели, а часть --- настройке
порога тревоги. При ограниченном объёме памяти увеличение одной части
уменьшает объём другой. Поэтому возникает задача выбора
их соотношения с учётом изменений трафика и возможного возвращения прежних
режимов. Цель работы состоит в разработке и экспериментальной оценке способов
такого распределения, направленных на снижение доли ложных тревог с учётом
сохранения обнаружения атак.

Исследование включает сравнительный обзор подходов к адаптации сетевых
детекторов, сохранению исторических данных и калибровке порогов, а также
формулировку задачи при общем ограничении на память. На этой основе будут
разработаны и реализованы способы определения долей наблюдений, сохраняемых
для дообучения и калибровки. Их оценка будет проводиться на воспроизводимых
последовательностях сетевых данных, содержащих нормальный трафик и атаки,
с изменением и повторным появлением нормальных режимов.

Предложенные способы планируется сравнить с вариантами, использующими
фиксированные доли обучающей и калибровочной памяти. Во всех сравниваемых
вариантах будет использоваться одинаковый общий бюджет памяти. Основными
показателями станут доля ложных тревог при возвращении нормальных режимов
и доля обнаруженных атак. Такое сравнение позволит оценить, помогает ли
выбранное распределение памяти сохранять качество детектора и не достигается
ли уменьшение ложных тревог за счёт пропуска атак.

Предполагаемыми результатами работы являются программная реализация
исследуемых способов, воспроизводимые сценарии экспериментальной проверки
и сравнительная оценка их эффективности. Ожидается определить, как соотношение
обучающей и калибровочной частей памяти влияет на поведение детектора,
и выявить условия полезности и ограничения каждого способа. Практическая
ценность результатов заключается в обосновании выбора распределения памяти
для сетевых детекторов, которые должны адаптироваться к текущему трафику
и учитывать возвращение ранее наблюдавшегося нормального поведения.

\abstractdocumentend
.
% При самостоятельной компиляции используется оформление предпросмотра.
\ifdefined\maketitle
    \let\abstractdocumentend\relax
\else
    \documentclass[12pt,a4paper]{article}
    \usepackage{fontspec}
    \setmainfont{cmunrm.otf}[BoldFont=cmunbx.otf,ItalicFont=cmunti.otf,BoldItalicFont=cmunbi.otf]
    \usepackage{polyglossia}
    \setdefaultlanguage{russian}
    \usepackage[top=20mm,bottom=20mm,left=30mm,right=15mm]{geometry}
    \usepackage{setspace}
    \onehalfspacing
    \setlength{\emergencystretch}{2em}
    \def\abstractdocumentend{\end{document}}
    \begin{document}
\fi

\section*{Аннотация}

Работа посвящена распределению ограниченной памяти между дообучением модели
и калибровкой порога сетевого детектора в условиях возвратного дрейфа трафика. Под возвратным дрейфом понимается повторное появление
ранее наблюдавшихся режимов сетевой активности после периода их отсутствия.
Например, дневная активность пользователей в сети организации
сменяется снижением нагрузки ночью. После адаптации к ночному
трафику детектор может ошибочно выдавать тревоги при возвращении
обычной дневной активности на следующее утро.
При непрерывном обновлении детектора такие ошибки могут возникать как
вследствие ухудшения распознавания прежних данных, так и из-за изменения
численных оценок аномальности, которым перестаёт соответствовать установленный
порог.

Для сохранения информации о прошлом используются исторические наблюдения.
Часть из них может служить повторному обучению модели, а часть --- настройке
порога тревоги. При ограниченном объёме памяти увеличение одной части
уменьшает объём другой. Поэтому возникает задача выбора
их соотношения с учётом изменений трафика и возможного возвращения прежних
режимов. Цель работы состоит в разработке и экспериментальной оценке способов
такого распределения, направленных на снижение доли ложных тревог с учётом
сохранения обнаружения атак.

Исследование включает сравнительный обзор подходов к адаптации сетевых
детекторов, сохранению исторических данных и калибровке порогов, а также
формулировку задачи при общем ограничении на память. На этой основе будут
разработаны и реализованы способы определения долей наблюдений, сохраняемых
для дообучения и калибровки. Их оценка будет проводиться на воспроизводимых
последовательностях сетевых данных, содержащих нормальный трафик и атаки,
с изменением и повторным появлением нормальных режимов.

Предложенные способы планируется сравнить с вариантами, использующими
фиксированные доли обучающей и калибровочной памяти. Во всех сравниваемых
вариантах будет использоваться одинаковый общий бюджет памяти. Основными
показателями станут доля ложных тревог при возвращении нормальных режимов
и доля обнаруженных атак. Такое сравнение позволит оценить, помогает ли
выбранное распределение памяти сохранять качество детектора и не достигается
ли уменьшение ложных тревог за счёт пропуска атак.

Предполагаемыми результатами работы являются программная реализация
исследуемых способов, воспроизводимые сценарии экспериментальной проверки
и сравнительная оценка их эффективности. Ожидается определить, как соотношение
обучающей и калибровочной частей памяти влияет на поведение детектора,
и выявить условия полезности и ограничения каждого способа. Практическая
ценность результатов заключается в обосновании выбора распределения памяти
для сетевых детекторов, которые должны адаптироваться к текущему трафику
и учитывать возвращение ранее наблюдавшегося нормального поведения.

\abstractdocumentend
.
% При самостоятельной компиляции используется оформление предпросмотра.
\ifdefined\maketitle
    \let\abstractdocumentend\relax
\else
    \documentclass[12pt,a4paper]{article}
    \usepackage{fontspec}
    \setmainfont{cmunrm.otf}[BoldFont=cmunbx.otf,ItalicFont=cmunti.otf,BoldItalicFont=cmunbi.otf]
    \usepackage{polyglossia}
    \setdefaultlanguage{russian}
    \usepackage[top=20mm,bottom=20mm,left=30mm,right=15mm]{geometry}
    \usepackage{setspace}
    \onehalfspacing
    \setlength{\emergencystretch}{2em}
    \def\abstractdocumentend{\end{document}}
    \begin{document}
\fi

\section*{Аннотация}

Работа посвящена распределению ограниченной памяти между дообучением модели
и калибровкой порога сетевого детектора в условиях возвратного дрейфа трафика. Под возвратным дрейфом понимается повторное появление
ранее наблюдавшихся режимов сетевой активности после периода их отсутствия.
Например, дневная активность пользователей в сети организации
сменяется снижением нагрузки ночью. После адаптации к ночному
трафику детектор может ошибочно выдавать тревоги при возвращении
обычной дневной активности на следующее утро.
При непрерывном обновлении детектора такие ошибки могут возникать как
вследствие ухудшения распознавания прежних данных, так и из-за изменения
численных оценок аномальности, которым перестаёт соответствовать установленный
порог.

Для сохранения информации о прошлом используются исторические наблюдения.
Часть из них может служить повторному обучению модели, а часть --- настройке
порога тревоги. При ограниченном объёме памяти увеличение одной части
уменьшает объём другой. Поэтому возникает задача выбора
их соотношения с учётом изменений трафика и возможного возвращения прежних
режимов. Цель работы состоит в разработке и экспериментальной оценке способов
такого распределения, направленных на снижение доли ложных тревог с учётом
сохранения обнаружения атак.

Исследование включает сравнительный обзор подходов к адаптации сетевых
детекторов, сохранению исторических данных и калибровке порогов, а также
формулировку задачи при общем ограничении на память. На этой основе будут
разработаны и реализованы способы определения долей наблюдений, сохраняемых
для дообучения и калибровки. Их оценка будет проводиться на воспроизводимых
последовательностях сетевых данных, содержащих нормальный трафик и атаки,
с изменением и повторным появлением нормальных режимов.

Предложенные способы планируется сравнить с вариантами, использующими
фиксированные доли обучающей и калибровочной памяти. Во всех сравниваемых
вариантах будет использоваться одинаковый общий бюджет памяти. Основными
показателями станут доля ложных тревог при возвращении нормальных режимов
и доля обнаруженных атак. Такое сравнение позволит оценить, помогает ли
выбранное распределение памяти сохранять качество детектора и не достигается
ли уменьшение ложных тревог за счёт пропуска атак.

Предполагаемыми результатами работы являются программная реализация
исследуемых способов, воспроизводимые сценарии экспериментальной проверки
и сравнительная оценка их эффективности. Ожидается определить, как соотношение
обучающей и калибровочной частей памяти влияет на поведение детектора,
и выявить условия полезности и ограничения каждого способа. Практическая
ценность результатов заключается в обосновании выбора распределения памяти
для сетевых детекторов, которые должны адаптироваться к текущему трафику
и учитывать возвращение ранее наблюдавшегося нормального поведения.

\abstractdocumentend
.
% При самостоятельной компиляции используется оформление предпросмотра.
\ifdefined\maketitle
    \let\abstractdocumentend\relax
\else
    \documentclass[12pt,a4paper]{article}
    \usepackage{fontspec}
    \setmainfont{cmunrm.otf}[BoldFont=cmunbx.otf,ItalicFont=cmunti.otf,BoldItalicFont=cmunbi.otf]
    \usepackage{polyglossia}
    \setdefaultlanguage{russian}
    \usepackage[top=20mm,bottom=20mm,left=30mm,right=15mm]{geometry}
    \usepackage{setspace}
    \onehalfspacing
    \setlength{\emergencystretch}{2em}
    \def\abstractdocumentend{\end{document}}
    \begin{document}
\fi

\section*{Аннотация}

Работа посвящена распределению ограниченной памяти между дообучением модели
и калибровкой порога сетевого детектора в условиях возвратного дрейфа трафика. Под возвратным дрейфом понимается повторное появление
ранее наблюдавшихся режимов сетевой активности после периода их отсутствия.
Например, дневная активность пользователей в сети организации
сменяется снижением нагрузки ночью. После адаптации к ночному
трафику детектор может ошибочно выдавать тревоги при возвращении
обычной дневной активности на следующее утро.
При непрерывном обновлении детектора такие ошибки могут возникать как
вследствие ухудшения распознавания прежних данных, так и из-за изменения
численных оценок аномальности, которым перестаёт соответствовать установленный
порог.

Для сохранения информации о прошлом используются исторические наблюдения.
Часть из них может служить повторному обучению модели, а часть --- настройке
порога тревоги. При ограниченном объёме памяти увеличение одной части
уменьшает объём другой. Поэтому возникает задача выбора
их соотношения с учётом изменений трафика и возможного возвращения прежних
режимов. Цель работы состоит в разработке и экспериментальной оценке способов
такого распределения, направленных на снижение доли ложных тревог с учётом
сохранения обнаружения атак.

Исследование включает сравнительный обзор подходов к адаптации сетевых
детекторов, сохранению исторических данных и калибровке порогов, а также
формулировку задачи при общем ограничении на память. На этой основе будут
разработаны и реализованы способы определения долей наблюдений, сохраняемых
для дообучения и калибровки. Их оценка будет проводиться на воспроизводимых
последовательностях сетевых данных, содержащих нормальный трафик и атаки,
с изменением и повторным появлением нормальных режимов.

Предложенные способы планируется сравнить с вариантами, использующими
фиксированные доли обучающей и калибровочной памяти. Во всех сравниваемых
вариантах будет использоваться одинаковый общий бюджет памяти. Основными
показателями станут доля ложных тревог при возвращении нормальных режимов
и доля обнаруженных атак. Такое сравнение позволит оценить, помогает ли
выбранное распределение памяти сохранять качество детектора и не достигается
ли уменьшение ложных тревог за счёт пропуска атак.

Предполагаемыми результатами работы являются программная реализация
исследуемых способов, воспроизводимые сценарии экспериментальной проверки
и сравнительная оценка их эффективности. Ожидается определить, как соотношение
обучающей и калибровочной частей памяти влияет на поведение детектора,
и выявить условия полезности и ограничения каждого способа. Практическая
ценность результатов заключается в обосновании выбора распределения памяти
для сетевых детекторов, которые должны адаптироваться к текущему трафику
и учитывать возвращение ранее наблюдавшегося нормального поведения.

\abstractdocumentend
